Теперь перейдём к операции умножения. Стандартный алгоритм умножения
$n$-битовых чисел $a$ и $b$ «в столбик» требует сложения не более чем $n$ двоичных
чисел, каждое из которых есть сдвиг $b$ на $j$ битов, если $a_j=1$. Если эти числа сложить при помощи двоичного дерева, получится схема размера $O(n^2)$ и глубины $O(\log^2 n)$: дерево имеет глубину $O(\log n)$, и глубина на каждом уровне тоже $O(\log n)$. Это даёт принадлежность умножения только к $\mathsf{}sf{NC}^2$. Покажем, что оно также лежит в $\mathsf{NC}^1$.

\begin{theorem} Функцию умножения можно вычислить схемой размера $O(n^2)$ и глубины $O(\log n)$.  \end{theorem}

\begin{proof}

Мы построим схему констанстной глубины, которая на вход получает три числа $x, y, z$ и возвращает два числа $u,v$, такие что $u + v = x + y + z$. После применения этой схемы $\log{\frac{2}{3} n} = O(\log n)$ раз останется 2 числа, которые сложим обычным образом. Таким образом, общая глубина составит те же $O(\log n)$.

Осталось построить такую схему. Идея проста: сложим $x,y,z$ без переносов, например в десятичной системе. В каждом разряде будут числа $0,1,2$ или $3$. Далее представим эти числа в двухбитовой двоичной записи: $00, 01, 10, 11$. Младшие биты всех этих чисел образуют число $u$, а старшие (со сдвигом на 1 влево) - число $v$. 

Рассмотрим для примера числа $x=1011, y=1101, z=1100$. Тогда $x+y+z = 3212 = 11100110$. Тогда $u = 1010$, а $v = 11010$ (обратите внимание, что $v$ дополнительно сдвинут на один влево). 

Глубина будет константной, ведь применяется одна и та же схема. Корректность конструкции почти очевидна: в число $u$ собраны все результаты сложения без переносов, а в число $v$ ~--- все биты переносов. \end{proof}